\section{问题四：考虑收缩的烘干时长}\label{sec:q4}

\subsection{求解}

附件 2 显示，半径在前 $\SI{6}{h}$ 由 $\SI{2}{cm}$ 降到 $\SI{1.374}{cm}$，$\SI{12}{h}$ 降到 $\SI{1.248}{cm}$，此后缓慢趋于 $\SI{1.2}{cm}$；收缩主要发生在快速失水的前期，与平均含水率的下降同步（图~\ref{fig:q4}(c)）。按第~\ref{sec:ref}~节的参考坐标，以附录 4 的物性和附件 2 的半径联立求解式~\eqref{eq:ref}、\eqref{eq:refbc}，取 $N=800$。result4.xlsx 以固定位置 $0,0.1,\dots,\SI{1.9}{cm}$ 和「药材表面」为列，位于药材之外的位置留空。

同样以 $\Cmax$ 降到 0.15 为判据（最大值仍在中心），得到烘干时长
\begin{equation}
\tstar=\SI{51.0920}{h}\quad(\SI{183931.30}{s})
\end{equation}
表~\ref{tab:t6} 给出每隔 $\SI{6}{h}$ 的含水率与当时的半径。预热平衡阶段在 $\SI{2.83}{h}$ 结束，完成了全部失水量的 35.6\%；干燥曲线同样符合双指数规律（图~\ref{fig:drying}）。

\begin{table}[H]\centering\small
\caption{药材烘干过程的水分浓度（\si{kg/kg}）}\label{tab:t6}
\begin{tabular}{cccccc}
\toprule
 & & \multicolumn{4}{c}{到药材中心的距离/cm} \\
\cmidrule(lr){3-6}
时间/h & 半径/cm & 0 & 0.5 & 1 & 药材表面 \\
\midrule
6 & 1.3740 & 1.7196 & 1.5378 & 1.0226 & 0.4207 \\
12 & 1.2480 & 0.7377 & 0.6546 & 0.4083 & 0.1669 \\
18 & 1.2140 & 0.4088 & 0.3687 & 0.2397 & 0.0892 \\
24 & 1.2040 & 0.2851 & 0.2611 & 0.1790 & 0.0673 \\
30 & 1.2010 & 0.2263 & 0.2094 & 0.1493 & 0.0595 \\
36 & 1.2000 & 0.1928 & 0.1797 & 0.1317 & 0.0560 \\
42 & 1.2000 & 0.1712 & 0.1604 & 0.1200 & 0.0541 \\
48 & 1.2000 & 0.1561 & 0.1468 & 0.1116 & 0.0530 \\
51.0920（烘干结束） & 1.2000 & 0.1500 & 0.1413 & 0.1081 & 0.0526 \\
\bottomrule
\end{tabular}

\end{table}

\begin{figure}[H]\centering
\includegraphics[width=\textwidth]{fig_q4.pdf}
\caption{问题四的结果。(a) 物理坐标下的含水率场，红线为表面 $R(t)$，橙色虚线为含水率 0.15 的等值线，右侧白色为药材之外；(b) 固定位置与表面的含水率；(c) 附件 2 的半径与计算的平均含水率；(d) 物性与收缩对烘干时长的分别影响（$N=400$）。}
\label{fig:q4}
\end{figure}

\subsection{物性与收缩的影响}

与问题三相比，问题四同时改变了物性和几何。为区分二者的作用，在同一网格（$N=400$）上计算三种情形（图~\ref{fig:q4}(d)）：附录 3 物性、固定半径，$\SI{57.47}{h}$；附录 4 物性、固定半径，$\SI{129.85}{h}$；附录 4 物性、按附件 2 收缩，$\SI{51.09}{h}$。

$\SI{50}{\celsius}$ 时附录 4 的扩散系数只有附录 3 的 0.19 倍（初始含水率）到 0.48 倍（含水率 0.15），若半径不变，烘干时长将延长到问题三的 2.3 倍。收缩使半径降到 $\SI{1.2}{cm}$，扩散距离缩短 40\%，按扩散时间与半径平方成正比估计约缩短到 36\%；两者相抵后，烘干时长反而比问题三短 $\SI{6.38}{h}$。对于会明显收缩的药材，忽略收缩算出的烘干时长是考虑收缩时的 2.5 倍。

把附件 2 的半径与计算的平均含水率逐时对应，得到半径与平均含水率的关系 $R=g(\Cbar)$。以它代替 $R(t)$ 重新求解，烘干时长为 $\SI{51.0917}{h}$，与按时间给定半径的结果仅差 $\SI{0.0004}{h}$。该关系把收缩与失水联系起来，可用于第~\ref{sec:process}~节改变工况后的计算。
